\documentclass[11pt]{article}
\usepackage[a4paper,margin=1in]{geometry}
\usepackage{fontspec}
\usepackage[boldfont=false]{xeCJK}
\setCJKmainfont{FandolSong-Regular.otf}
\usepackage{amsmath}
\usepackage{amssymb}
\usepackage{array}
\usepackage{booktabs}
\usepackage{graphicx}
\usepackage{adjustbox}
\usepackage{enumitem}
\setlist[itemize]{topsep=2pt,partopsep=0pt,parsep=0pt,itemsep=2pt,leftmargin=1.4em}
\setlist[enumerate]{topsep=2pt,partopsep=0pt,parsep=0pt,itemsep=2pt,leftmargin=1.6em}
\usepackage{parskip}
\usepackage{fvextra}
\fvset{breaklines=true,breakanywhere=true,fontsize=\small}
\setlength{\parindent}{0pt}

\begin{document}

\begin{center}
  {\LARGE\bfseries Energy and respiration}\\[0.35em]
  {\large A Level Biology}
\end{center}
\vspace{0.6em}

\section*{Why living things need energy}

Cells need a steady supply of \textbf{energy} 能量, which they get from \textbf{respiration} 呼吸作用. Energy is needed for:

\begin{itemize}
\item \textbf{active transport} 主动运输 (moving substances against a gradient),

\item movement (for example muscle contraction),

\item \textbf{anabolic} 合成代谢 reactions — the building of large molecules, such as in DNA replication and protein synthesis.

\end{itemize}

\section*{ATP — the energy currency}

ATP is the molecule that carries energy to where it is needed. It is made from ADP and a phosphate group; when it loses that phosphate again, it releases a small, usable burst of energy. ATP suits this job well, so we call it the \textbf{universal energy currency}:

\begin{itemize}
\item it releases energy quickly, in small amounts that match a cell's needs.

\item it is easily made and re-made, again and again.

\item it is small and soluble, so it moves easily around the cell.

\end{itemize}

ATP is made in two ways: by direct transfer of a phosphate group in \textbf{phosphorylation} 磷酸化 reactions, and by \textbf{chemiosmosis} 化学渗透 across the membranes of \textbf{mitochondria} 线粒体 and chloroplasts.

\section*{Respiratory substrates and RQ}

A \textbf{respiratory substrate} 呼吸底物 is a molecule that is broken down to release energy. Per gram, \textbf{lipids} release the most energy (they have the most hydrogen), proteins are next, and carbohydrates the least.

The \textbf{respiratory quotient} 呼吸商 (RQ) compares the gases exchanged:

\[
\text{RQ} = \frac{\text{molecules of carbon dioxide produced}}{\text{molecules of oxygen taken in}}
\]

You can work out the RQ from a respiration equation. Carbohydrates give an RQ of about 1.0, lipids about 0.7 and proteins about 0.9.

A \textbf{respirometer} 呼吸计 measures the \textbf{oxygen} 氧气 taken in by living things, such as \textbf{germinating} 萌发 seeds or small \textbf{invertebrates} 无脊椎动物, and is used to find their RQ.

\section*{Aerobic respiration: the four stages}

\textbf{Aerobic} 有氧 respiration (with oxygen) has four stages, each in a set place in the cell:

\begin{center}
\adjustbox{max width=\linewidth}{%
\begin{tabular}{ll}
\toprule
\textbf{Stage} & \textbf{Where it happens} \\
\midrule
\textbf{glycolysis} 糖酵解 & the \textbf{cytoplasm} 细胞质 \\
\textbf{link reaction} 连接反应 & the \textbf{matrix} 基质 of the mitochondria \\
\textbf{Krebs cycle} 克雷布斯循环 & the matrix of the mitochondria \\
\textbf{oxidative phosphorylation} 氧化磷酸化 & the inner membrane of the mitochondria \\
\bottomrule
\end{tabular}}
\end{center}

\subsection*{Glycolysis}

\textbf{Glucose} 葡萄糖 (6 carbons) is first phosphorylated, using 2 ATP, to form fructose bisphosphate (6C). This 6C molecule is split into two triose phosphate molecules (3C each). These are then \textbf{oxidised} 氧化 to \textbf{pyruvate} 丙酮酸 (3C). Glycolysis makes a net gain of 2 ATP and some \textbf{reduced} 还原 NAD (NAD is a \textbf{coenzyme} 辅酶, a helper molecule).

\subsection*{The link reaction}

When oxygen is available, pyruvate enters the mitochondria. There each pyruvate loses a \textbf{carbon dioxide} 二氧化碳 and is turned into a 2-carbon \textbf{acetyl} 乙酰基 group. This group is carried by \textbf{coenzyme A} 辅酶A to form acetyl coenzyme A. Some carbon dioxide is released and NAD is reduced.

\subsection*{The Krebs cycle}

The 2C acetyl group joins a 4-carbon molecule, \textbf{oxaloacetate} 草酰乙酸, to make a 6-carbon molecule, \textbf{citrate} 柠檬酸. Citrate is then changed back to oxaloacetate in a series of small steps, ready to accept the next acetyl group. During these steps:

\begin{itemize}
\item \textbf{decarboxylation} 脱羧 removes carbon as carbon dioxide.

\item \textbf{dehydrogenation} 脱氢 removes hydrogen, which \textbf{reduces} the coenzymes NAD and FAD.

\end{itemize}

The reduced NAD and FAD then carry the hydrogen to the carriers in the inner mitochondrial membrane.

\subsection*{Oxidative phosphorylation}

This stage makes most of the ATP:

\begin{enumerate}
\item the hydrogen atoms split into \textbf{protons} 质子 and energetic \textbf{electrons} 电子.

\item the electrons pass along the \textbf{electron transport chain} 电子传递链, releasing energy as they go.

\item this energy is used to pump protons across the inner membrane.

\item the protons flow back into the matrix through a channel called \textbf{ATP synthase} 合酶. This flow provides the energy to make ATP (this is chemiosmosis).

\item \textbf{oxygen} is the final electron acceptor: it joins with electrons and protons to form water.

\end{enumerate}

\subsection*{The structure of mitochondria}

The inner membrane is folded into \textbf{cristae} 嵴, giving a large surface for the electron transport chain and ATP synthase. The matrix inside holds the substances and helpers for the link reaction and the Krebs cycle.

\section*{Anaerobic respiration}

When there is no oxygen, only glycolysis can run. To keep glycolysis going, the cell must use up the reduced NAD. This happens by \textbf{fermentation} 发酵:

\begin{itemize}
\item in mammals, pyruvate is turned into \textbf{lactate} 乳酸 (lactate fermentation). The lactate is later broken down when oxygen returns.

\item in \textbf{yeast} 酵母, pyruvate is turned into \textbf{ethanol} 乙醇 and carbon dioxide (ethanol fermentation).

\end{itemize}

\textbf{Anaerobic} 无氧 respiration gives far less energy than aerobic respiration. Aerobic respiration also runs the Krebs cycle and oxidative phosphorylation, which release a lot more ATP, while anaerobic respiration gains only the small amount from glycolysis.

\section*{Rice and waterlogged roots}

Rice can grow with its roots under water, where there is little oxygen. It is adapted in three ways: it develops \textbf{aerenchyma} 通气组织 (air-filled spaces) in the roots to carry air down; the roots use ethanol fermentation to keep making some ATP; and the stems grow faster to reach the air above the water.

\section*{Investigating the rate of respiration}

A redox \textbf{indicator} 指示剂 such as DCPIP or methylene blue loses its colour when it gains hydrogen from respiring cells. The faster the colour is lost, the faster the yeast is respiring, so you can test the effect of temperature or substrate concentration. A respirometer can also be used to measure how temperature changes the rate of oxygen uptake.


\end{document}
